%HOMEWORK template for M 325K
\documentclass{article}
\headheight=7pt         \topmargin=14pt
\textheight=574pt       \textwidth=445pt
\oddsidemargin=18pt     \evensidemargin=18pt

%Begin package import

\usepackage{amsmath, amssymb, amsthm}
\usepackage{mathtools}
\usepackage{pdfpages}
\usepackage{tikz-cd}

%End package import
%Begin custom commands

\newcommand*{\T}{\mathcal{T}}
\newcommand*{\B}{\mathcal{B}}
\newcommand*{\U}{\mathcal{U}}
\newcommand*{\cl}{\mathrm{cl}}
\newcommand*{\subeq}{\subseteq}
\newcommand*{\empset}{\emptyset}
\newcommand{\C}{\mathbb{C}}
\newcommand{\R}{\mathbb{R}}
\newcommand{\Q}{\mathbb{Q}}
\newcommand{\Z}{\mathbb{Z}}
\newcommand{\N}{\mathbb{N}}
\newcommand{\F}{\mathbb{F}}
\newcommand{\abs}[1]{\left\lvert #1\right\rvert}
\newcommand{\RM}[1]{\mathrm{#1}}
\newcommand{\BF}[1]{\textbf{#1}}
\newcommand{\e}{\epsilon}
\newcommand{\ceil}[1]{\lceil{#1}\rceil}
\newcommand{\floor}[1]{\lfloor{#1}\rfloor}
\newcommand{\rarr}[0]{\rightarrow}
\newcommand{\IM}[1]{\text{Im}(#1)}
\newcommand{\Hom}[0]{\text{Hom}}
\newcommand{\Pow}[1]{\text{Pow}(#1)}
\newcommand{\angles}[1]{\langle #1 \rangle}
\newcommand{\Ext}[0]{\text{Ext}}
\newcommand{\Tor}[0]{\text{Tor}}
\newcommand{\rank}[0]{\text{rank}}
\newcommand{\CP}[1]{\C P^{#1}}
\newcommand{\x}[0]{\times}




%End custom commands
%Begin custom theorems

\newtheorem{innercustomgeneric}{\customgenericname}
\providecommand{\customgenericname}{}
\newcommand{\newcustomtheorem}[2]{%
\newenvironment{#1}[1]
{%
\renewcommand\customgenericname{#2}%
\renewcommand\theinnercustomgeneric{##1}%
\innercustomgeneric%
}
{\endinnercustomgeneric}
}

\newcustomtheorem{THRM}{Theorem}
\newcustomtheorem{LEM}{Lemma}
\newcustomtheorem{EX}{Exercise}
\newcustomtheorem{COR}{Corollary}
\newcustomtheorem{PROB}{Problem}
\newcustomtheorem{claim}{Claim}

\newenvironment{solution}
{\renewcommand\qedsymbol{$\blacksquare$}\begin{proof}[Solution]}
{\end{proof}}

\newenvironment{definition}
{\renewcommand\qedsymbol{$\blacksquare$}\begin{proof}[Definition]}
{\end{proof}}

%End custom theorems
\begin{document}

\title{Hatcher Chapter 3: Cohomology}
\author{Arham Lodha}%put your name here
\date{May 13, 2025 - }
\maketitle

\section{Universal Coefficent Theorem for Cohomology}
\section{Cup Product}



\subsection{Problems}
\begin{PROB}{3.2.5}\label{Problem 3.2.5.}
	Show the ring $H^*(\R P^\infty, \Z_{2k}) \cong \Z_{2k} [\alpha, \beta] / (2\alpha, 2\beta, \alpha^2 - k \beta)$
\end{PROB}
\begin{proof}
	By Theorem 3.19, $H^*(\R P^{\infty}, \Z_{2}) = \Z[a]$.
	Use standard cellular structure and standard cellular cohomology to see that

	\[H^0(\R P^n, \Z_{2k}) = \Z_{2 k}\], \[H^n(\R P^n, \Z_{2 k}) = \Z_2, \forall n \geq 1\]


	Let $r: \Z_{2 k} \to \Z_2$ be the Coefficent maps reduction mod $2$ ie $[x]_{2k} \to [x]_2$. It induces a surjective ring map, $r^*: H^*(\R P^\infty, \Z_{2 k}) \to H^*(\R P^{\infty}, \Z_2)$

	Let $\alpha \in H^1(\R P^{\infty}, \Z_{2k})$ where $r^*(\alpha) = a \in H^1(\R P^{\infty}, \Z_2) = \Z[a]$, since $\alpha$ maps to the generator (only non 0 element) it must also be nonzero. Since $H^1(\R P^{\infty}, \Z_{2k}) = \Z_2$, $\alpha$ generates $H^1$.

	Let $\beta \in H^2(\R P^{\infty}, \Z_{2k})$ where $r^*(\beta) = a^2$ and $\beta$ generates $H^2(\R P^{\infty}, \Z_{2k})$.

	Note since $H^n = \Z_2$, $2\alpha = 0$ and $2\beta = 0$. So the $(2\alpha, 2\beta)$ is forced on the additive structure.

	Now lets see how $\alpha^2$ and $\beta$ relate.

	Take the short exact sequence

	\[0 \to \Z_k \xhookrightarrow{i} \Z_{2k} \to \Z_2 \to 0\]

	The Bockstein Homomorphims induced by the coefficient sequence:

	$B(a) = a^2 = k \beta$. So naturality gives $r^*(\alpha) = a^2 = r^*(k \beta) \implies \alpha - k \beta = 0$. So the ideal $(\alpha - k \beta)$ gets added.


	Thus $H^*(\R P^\infty, \Z_{2k}) = \Z_{2 k}[\alpha, \beta] / (2\alpha, 2\beta, \alpha^2 - k \beta)$, where $\abs{\alpha} = 1$ and $\abs{\beta} = 2$.
\end{proof}

\bigskip

\begin{PROB}{3.2.6}\label{Problem 3.2.6.}
	Use cup products to compute the map $H^*(\C P^n, \Z) \to H^*(\C P^n, \Z)$ induced by the map $\C P^n \to \C P^n$ that this the quotient of the map $\C^{n + 1} \to \C^{n + 1}$ raising each coordinate to the $d$th power.
\end{PROB}
\begin{proof}
	Let $P^n = \C P^n$ and all coefficients are $\Z$ unless otherwise stated.

	\begin{enumerate}
		\item For $n = 1$: $\C P^1 = S^2$, by Riemann Sphere. Thus the map $H^*(S^2) \to H^*(S^2)$ is just the $\times d$ map on $H^2$ and 0 map everywhere else.
		\item For $n > 1$: $i: \C P^1 \to \C P^n$ be the natural inclusion.

            % https://q.uiver.app/#q=WzAsNCxbMCwwLCJIXjIoUF5uKSJdLFsxLDAsIkheMihQXm4pIl0sWzAsMSwiSF4xKFBeMSkiXSxbMSwxLCJIXjEoUF4xKSJdLFsyLDMsImQiXSxbMCwxLCJkIl0sWzAsMiwiaV4qIiwyXSxbMSwzLCJpXioiXV0=
            \[\begin{tikzcd}
                    {H^2(P^n)} & {H^2(P^n)} \\
                    {H^1(P^1)} & {H^1(P^1)}
                    \arrow["d", from=1-1, to=1-2]
                    \arrow["{i^*}"', from=1-1, to=2-1]
                    \arrow["{i^*}", from=1-2, to=2-2]
                    \arrow["d", from=2-1, to=2-2]
                \end{tikzcd}\]

            $i^*$ is a isomorphism. So top d must be $d \times$.
	\end{enumerate}
\end{proof}

\bigskip

\begin{PROB}{3.2.7}\label{Problem 3.2.7.}
	Use cup products to show $\R P^3$ is not homotopy equivalent to $\R P^2 \vee S^3$
\end{PROB}
\begin{proof}
	\begin{enumerate}
		\item \textbf{Without cup products}: Universal cover of $P^3$ is $S^3$ and universal cover of $P^2 \vee S^3 = S^3 \vee S^2 \vee S^3$. $H_3(S^3) = \Z$ and $H_3(S^3 \vee S^2 \vee S^3) = \Z \oplus \Z$. Thus universal covers are not homeomorphic. Thus the spaces are not homotopy equivalent.
		\item \textbf{With Cup products}: Unless explicitly stated $\Z_2$ coefficients are being used. By Theorem 3.19:

		      \[
			      H^*(P^3) := \Z_2[\alpha]/(\alpha^4)
		      \].

		      By Cellular Cohomology and 3.14:

		      \[H^*(P^2 \vee S^3) = \frac{\Z_2[\alpha]}{(\alpha^3)} \times \frac{\Z_2 [\beta]}{(\beta^2)}\]

		      where $\abs{\alpha} = 1$ and $\abs{\beta} = 3$. Every degree 1 class has the form, $(\lambda \alpha, 0) \implies (\lambda \alpha, 0)^3 = (\lambda^3 \alpha^3, 0 ) = (0, 0) = 0$. Thus no element of degree 1 has a nonzero cube. Thus both graded rings are not isomorphic. Thus both spaces are not homotopy equivalent.
	\end{enumerate}

\end{proof}

\bigskip

\begin{PROB}{3.2.8}\label{Problem 3.2.8.}
	Let $X = \C P^2$ with a cell $e^3$ attached by a map $S^2 = \C P^1 \subset \C P^2$ with degree $p$, and let $Y = M(\Z_p, 2) \vee S^4$. Thus X and Y have the same 3 skeleton but differ in the way
	their 4 cells are attached. Show that X and Y have isomorphic cohomology rings with $\Z$ coefficients but not with $\Z_p$ coefficients.
\end{PROB}
\begin{proof}
	\begin{enumerate}
		\item $\Z$ coefficients:
		      \begin{enumerate}
			      \item Cohomology of $X$ and $Y$:


			            \[\begin{tikzcd}
					            0 & \Z & 0 & \Z & \Z & \Z & 0
					            \arrow[from=1-1, to=1-2]
					            \arrow[from=1-2, to=1-3]
					            \arrow[from=1-3, to=1-4]
					            \arrow["{\times p}"', from=1-4, to=1-5]
					            \arrow["0"', from=1-5, to=1-6]
					            \arrow[from=1-6, to=1-7]
				            \end{tikzcd}\]

			            \[H^n(X) := \begin{cases*}
					            \Z   & if $n = 0, 4$            \\
					            \Z_p & if $n = 3$               \\
					            0    & if $n = 1, 2$ or $n > 4$
				            \end{cases*}\]

                    It makes sense that for degree reasons that the cup product structure is trivial. (Cup product of any two nonzero homology generators)
                            
		      \end{enumerate}

		\item $\Z_p$ coefficients

		      \begin{enumerate}
			      \item Cohomology of X:

			            % https://q.uiver.app/#q=WzAsNyxbMiwwLCIwIl0sWzEsMCwiXFxaIl0sWzAsMCwiMCJdLFszLDAsIlxcWiJdLFs0LDAsIlxcWiJdLFs1LDAsIlxcWiJdLFs2LDAsIjAiXSxbMiwxXSxbMSwwXSxbMCwzXSxbMyw0LCJcXHRpbWVzIHAiLDJdLFs0LDUsIjAiLDJdLFs1LDZdXQ==
			            \[\begin{tikzcd}
					            0 & \Z_p & 0 & \Z_p & \Z_p & \Z_p & 0
					            \arrow[from=1-1, to=1-2]
					            \arrow[from=1-2, to=1-3]
					            \arrow[from=1-3, to=1-4]
					            \arrow["{\times p}"', from=1-4, to=1-5]
					            \arrow["0"', from=1-5, to=1-6]
					            \arrow[from=1-6, to=1-7]
				            \end{tikzcd}\]

			            \[H^n(X) := \begin{cases*}
					            \Z_p & if $n = 0, 2, 3, 4$   \\
					            0    & if $n = 1$ or $n > 4$
				            \end{cases*}\]

                        The cup product is only nontrivial on the 2 chain. So the cup product structure reduces to the cup product structure of $\C P^2$. Let $\alpha \in H^2(X)$, be the generator by Theorem 3.19, $\alpha \smile \alpha$ is the generator of $H^4(X)$.  
                        
                \item Cohomology of $Y$
                
                % https://q.uiver.app/#q=WzAsNyxbMiwwLCIwIl0sWzEsMCwiXFxaIl0sWzAsMCwiMCJdLFszLDAsIlxcWiJdLFs0LDAsIlxcWiJdLFs1LDAsIlxcWiJdLFs2LDAsIjAiXSxbMiwxXSxbMSwwXSxbMCwzXSxbMyw0LCJcXHRpbWVzIHAiLDJdLFs0LDUsIjAiLDJdLFs1LDZdXQ==
			            \[\begin{tikzcd}
                            0 & \Z_p & 0 & \Z_p & \Z_p & \Z_p & 0
                            \arrow[from=1-1, to=1-2]
                            \arrow[from=1-2, to=1-3]
                            \arrow[from=1-3, to=1-4]
                            \arrow["{\times p}"', from=1-4, to=1-5]
                            \arrow["0"', from=1-5, to=1-6]
                            \arrow[from=1-6, to=1-7]
                        \end{tikzcd}\]

                    \[H^n(X) := \begin{cases*}
                            \Z_p & if $n = 0, 2, 3, 4$   \\
                            0    & if $n = 1$ or $n > 4$
                        \end{cases*}\]
                    
                    The reduced Cohomology ring is:

                    \[
                    \tilde{H}^*(Y) := \tilde{H}^*(M(\Z_p, 2)) \times \Z[\beta]/[\beta^2]
                    \]

                    for $\abs{\beta} = 4$. Here the generator of the $H^2$, cup product with itself will be trivial.   

		      \end{enumerate}

	\end{enumerate}
\end{proof}

\bigskip

\begin{PROB}{3.2.9}\label{Problem 3.2.9.}
    Show that if $H_n(X, \Z)$ are free for all $n$, then $H^*(X, \Z_p)$ and $H^*(X, \Z) \otimes \Z_p$ are isomorphic as rings, so in particular the ring structure with $\Z$ coefficients determines the ring structure with $\Z_p$ coefficients.      
\end{PROB}
\begin{proof}
    By Universal Coefficient Theorem for Homology (Theorem 3A.3), 
    
    \[H_n(X, \Z_p) = H_n(Z) \otimes \Z_p \oplus \Tor\left(H_{n - 1}(X), \Z_p\right) = H_n(\Z) \otimes \Z_p\]

    The last equality is true because $H_n(X)$ is free for all $n$, furthermore $H_n(\Z) \otimes \Z_p = \Z_p^{\rank H_n(\Z)}$ ie a free module over $\Z_p$.   

    By Universal Coefficient Theorem for Cohomology (Theorem 3.2):
    \[H^n(X, \Z_p) = \Ext(H_{n - 1}(X, \Z_p), \Z_p) \oplus \Hom(H_n(X, \Z_p), \Z_p) = \Hom(H_n(X, \Z_p), \Z_p) = \Z_p^{\rank H_n(\Z)} \]

    Since $H_{n - 1}(X, \Z_p)$ is free over $\Z_p$, Ext term disappears.
    
    \[H^n(X, \Z) = \Ext(H_{n - 1}(X, \Z), \Z) \oplus \Hom(H_n(X, \Z), \Z) = \Z^{\rank H_n(X)}\]

    Thus $H^n(X, \Z) \otimes \Z_p = \Z_p^{\rank H_n(X)} \cong H^n(X, \Z_p)$. Thus there is a graded ring isomorphism. 
    
    Let $\pi: \Z \to \Z_p$ standard mod map. Let $\rho_p: C^*(X, \Z) \to C^*(X, \Z_p)$ be the standard coefficient reduction map where $c \to \pi \circ c$. $\pi$ is a ring map. Let $\alpha \in C^p(X, \Z)$, $\beta \in C^{q}(X, \Z)$ and $\sigma$ be a $(p + q + 1)$ simplex, 
    
    \begin{align*}
        \rho_p(\alpha \smile \beta)(\sigma) &= (\pi \circ (\alpha \smile \beta))(\sigma)\\
        &= \pi(\alpha(\sigma|_{[v_0, \dots, v_p]}) \beta(\sigma|_{[v_{p + 1}, \dots, v_{p + q}]})) \\
        &= \pi(\alpha(\sigma|_{[v_0, \dots, v_p]})) \pi( \beta(\sigma|_{[v_{p + 1}, \dots, v_{p + q}]})) \\
        &= (\pi \circ \alpha)(\sigma|_{[v_0, \dots, v_p]}) (\pi \circ \beta)(\sigma|_{[v_{p + 1}, \dots, v_{p + q}]}) \\
        &= \rho_p(\alpha)(\sigma|_{[v_0, \dots, v_p]}) \rho_p(\beta)(\sigma|_{[v_{p + 1}, \dots, v_{p + q}]}) \\
        &= (\rho_p(\alpha) \smile \rho_p(\beta))(\sigma)
    \end{align*}

    Thus $\rho_p(\alpha \smile \beta) = \rho_p(\alpha) \smile \rho_p(\beta)$. 
    
    \[\rho_p^*: H^*(X, \Z) \to H^*(X, \Z_p)\]

    is a ring homomorphism. Let 
    
    \[\varphi_p = \rho_p^* \otimes 1_{\Z_p} : H^*(X, \Z) \otimes_\Z \Z_p \to H^*(X, \Z_p)\]

    $\varphi_p$ is the graded abelian group iso seen above. Tensoring with $\Z_p$ doesn't change the ring homomorphism, thus $\varphi_p$ is a ring isomorphism.    
\end{proof}

\bigskip

\begin{PROB}{3.2.10}\label{Problem 3.2.10.}
    Show that the cross product map $H^*(X, \Z) \otimes H^*(Y, \Z) \to H^*(X \times Y, \Z)$ is not a isomorphism if $X$ and $Y$ are infinite discrete sets.   
\end{PROB}
\begin{proof}
    Note $H^*(X, \Z) = \Pi_{x \in X} \Z$ similarly for $Y$. $H^*(X \times Y, \Z) = \Pi_{(x, y) \in X \times Y} \Z$. 
    
    \[\Pi_{x \in X} \Z \otimes \Pi_{y \in Y} \Z \subset \Pi_{(x, y) \in X \times Y} \Z\]

    where the inclusion is the natural one, 
    \begin{align*}
        (n_x)_{x \in X} \otimes (m_x)_{y \in Y} \to (n_x m_y)_{x \in X, y \in Y}
    \end{align*}'

    \textbf{UNFINISHED}
\end{proof}

\bigskip

\begin{PROB}{3.2.11}\label{Problem 3.2.11.}
    Using cup products show that every map $S^{m + n} \to S^{m} \times S^n$ induces the trivial homomorphism $H_{m + n}(S^{m + n}) \to H_{m + n}(S^m \times S^n)$, assuming $m > 0$ and $n > 0$.     
\end{PROB}
\begin{proof}
	It is known

	\[
	H^k(S^n) := \begin{cases*}
		\Z & if $k = 0, n$ \\
		0 & otherwise 
	\end{cases*}
	\]

	$\forall f: S^{m + n} \to S^m \times S^n$. By Proposition 3.10, $f$ induces 
	
	\[f^*: H^{*}(S^{m} \times S^{n}) \to H^*(S^{m + n})\]

	By Theorem 3.15, $H^{i}(S^n) \otimes_\Z H^{j}(S^m) \cong H^{i+j}(S^{n} \times S^m)$ and $p_1(a) \smile p_2(b)$ is a generator for $H^{i+j}(S^{n} \times S^m)$ where $a$ and $b$ are generators.      
	$\forall (a, b) \in H^{i}(S^{n}) \otimes_{\Z} H^{j}(S^{m})$. 

	\begin{align*}
		f^{i+j}(p^i_1(a) \smile p^j_2(b)) &= f^{i}(p^i_1(a)) \smile f^{j}(p^j_2(b)) \\
		&= (p_1 \circ f)^i(a) \smile (p_2 \circ f)^j(b)
	\end{align*} 

	we get the following commutative diagram for all $0 < i, j < n$ 

	% https://q.uiver.app/#q=WzAsNixbMiwwLCJIXntpK2p9KFNee20rbn0pIl0sWzIsMSwiSF57aStqfShTXnttfSBcXHRpbWVzIFNee259KSJdLFswLDAsIkheaShTXnttICsgbn0pIl0sWzAsMSwiSF5pKFNebSkiXSxbMSwxLCJIXmooU15uKSJdLFsxLDAsIkheaihTXnttICsgbn0pIl0sWzEsMCwiZl57aStqfSIsMl0sWzMsNCwiXFxvdGltZXNfXFxaIiwxLHsic3R5bGUiOnsiYm9keSI6eyJuYW1lIjoibm9uZSJ9LCJoZWFkIjp7Im5hbWUiOiJub25lIn19fV0sWzIsNSwiXFx0aW1lcyIsMSx7InN0eWxlIjp7ImJvZHkiOnsibmFtZSI6Im5vbmUifSwiaGVhZCI6eyJuYW1lIjoibm9uZSJ9fX1dLFszLDIsIihwXzEgXFxjaXJjIGYpXmkiXSxbNCw1LCIocF8yIFxcY2lyYyBmKV5qIl0sWzUsMCwiXFxzbWlsZSJdLFs0LDEsIlxcdGltZXMiLDJdXQ==
\[\begin{tikzcd}
	{H^i(S^{m + n})} & {H^j(S^{m + n})} & {H^{i+j}(S^{m+n})} \\
	{H^i(S^m)} & {H^j(S^n)} & {H^{i+j}(S^{m} \times S^{n})}
	\arrow["\times"{description}, draw=none, from=1-1, to=1-2]
	\arrow["\smile", from=1-2, to=1-3]
	\arrow["{(p_1 \circ f)^i}", from=2-1, to=1-1]
	\arrow["{\otimes_\Z}"{description}, draw=none, from=2-1, to=2-2]
	\arrow["{(p_2 \circ f)^j}", from=2-2, to=1-2]
	\arrow["\times"', from=2-2, to=2-3]
	\arrow["{f^{i+j}}"', from=2-3, to=1-3]
\end{tikzcd}\]
	
	Note $\times$ is a isomorphism. Furthermore note that $(p_1 \circ f)^i$ and $(p_2 \circ f)^j$ are trivial because $H^{i}(S^{m + n})$ and $H^j(S^{m + n})$ are 0. Thus $f^{i + j}$ must be 0. Thus $f^{m + n} = 0$.          


	By the Universal Coefficent Theorem for Cohomology we get the following: 

	% https://q.uiver.app/#q=WzAsMTAsWzIsMCwiSF57bStufShTXm0gXFx0aW1lcyBTXm4pIl0sWzEsMCwiXFxFeHQoSF97bStuIC0gMX0oU15tIFxcdGltZXMgU15uKSwgXFxaKSJdLFswLDAsIjAiXSxbMywwLCJcXEhvbShIX3ttK259KFNebSBcXHRpbWVzIFNebiksIFxcWikiXSxbNCwwLCIwIl0sWzAsMSwiMCJdLFsxLDEsIlxcRXh0KEhfe20rbiAtIDF9KFNee20rbn0pLCBcXFopIl0sWzIsMSwiSF57bStufShTXnttK259KSJdLFszLDEsIlxcSG9tKEhfe20rbn0oU157bStufSksIFxcWikiXSxbNCwxLCIwIl0sWzIsMV0sWzEsMF0sWzAsM10sWzMsNF0sWzEsNiwiKGZfKileKiJdLFs1LDZdLFs2LDddLFs3LDhdLFs4LDldLFszLDgsIihmXyopXioiXSxbMCw3LCJmXio9MCJdXQ==
\[\begin{tikzcd}
	0 & {\Ext(H_{m+n - 1}(S^m \times S^n), \Z)} & {H^{m+n}(S^m \times S^n)} & {\Hom(H_{m+n}(S^m \times S^n), \Z)} & 0 \\
	0 & {\Ext(H_{m+n - 1}(S^{m+n}), \Z)} & {H^{m+n}(S^{m+n})} & {\Hom(H_{m+n}(S^{m+n}), \Z)} & 0
	\arrow[from=1-1, to=1-2]
	\arrow[from=1-2, to=1-3]
	\arrow["{(f_*)^*}", from=1-2, to=2-2]
	\arrow[from=1-3, to=1-4]
	\arrow["{f^*=0}", from=1-3, to=2-3]
	\arrow[from=1-4, to=1-5]
	\arrow["{(f_*)^*}", from=1-4, to=2-4]
	\arrow[from=2-1, to=2-2]
	\arrow[from=2-2, to=2-3]
	\arrow[from=2-3, to=2-4]
	\arrow[from=2-4, to=2-5]
\end{tikzcd}\]

	all the $\Ext$s are 0, because all homology groups are free. Thus 3 arrows in each short exact sequence is a isomorphism. Thus if $f^* = 0 \implies (f_*)^* = 0 \implies f_* = 0$.   
\end{proof}

\bigskip

\begin{PROB}{3.2.12}\label{Problem 3.2.12.}
	Show that the spaces $(S^1 \x \CP{\infty}) / (S^1 \x \left\{ x_0 \right\})$ and $S^3 \x \CP{\infty}$ have isomorphic cohomology rings with $\Z$ and any other coefficients.  
\end{PROB}
\begin{proof}
	\textbf{$\Z$ coefficients}: By Example 3.14 and Theorem 3.19, 

	\[H^*(S^3 \x \CP{\infty}) \cong \frac{\Z[\alpha]}{\alpha^2} \otimes_{\Z} \Z[\beta] \cong \frac{\Z[\alpha, \beta]}{\alpha^2}\]

	where $\abs{\alpha} = 3$ and $\abs{\beta} = 2$. 
	
	By Excision , 
	\[\widetilde{H}^k((S^1 \x \CP{\infty}) / (S^1 \x \left\{ x_0 \right\})) = H^k(S^1 \x \CP{\infty} , S^1 \x \left\{ x_0 \right\}) = H^k(S^1 \x \CP{\infty}, \empset \x \CP{\infty} \cup S^1 \x \left\{ x_0 \right\})\] 

	By Relative Kunneth's Formula (Theorem 3.18),
	

	\begin{align*}
		H^k(S^1 \x \CP{\infty}, \empset \x \CP{\infty} \cup S^1 \x \left\{ x_0 \right\}) &= \bigoplus_{n = 0}^k H^n(S^1, \empset) \otimes_{\Z} H^{k - n}(\CP{\infty}, \left\{ x_0 \right\}) \\
		&= \bigoplus_{n = 0}^k H^n(S^1) \otimes_{\Z} \widetilde{H}^{k - n}(\CP{\infty}) \\
		&= \Z \otimes_{\Z} \widetilde{H}^{k}(\CP{\infty}) \oplus \Z \otimes_{\Z} \widetilde{H}^{k - 1}(\CP{\infty})
	\end{align*}

	


	Thus 

	\[H^k(S^1 \x \CP{\infty}, \empset \x \CP{\infty} \cup S^1 \x \left\{ x_0 \right\}) = \begin{cases*}
		0 & if $k \leq 1$ \\
		\Z & otherwise
	\end{cases*}\]

	Thus 

	\[H^k((S^1 \x \CP{\infty}) / (S^1 \x \left\{ x_0 \right\})) = \begin{cases*}
		0 & if $k = 1$ \\
		\Z & otherwise 
	\end{cases*}\]

	Let $a \in H^1(S^1)$ and $b \in \widetilde{H}^2(\CP{\infty})$ be generators. Then by the Kunneth isomorphism

	\[H^{k}(S^1 \x \CP\infty, S^1 \x \left\{ x_0 \right\}) \cong (H^0(S^1) \otimes \widetilde{H}^k(\CP\infty)) \oplus (H^1(S^1) \otimes \widetilde{H}^{k-1}(\CP\infty))\]

	so we have 2 classes, $\alpha = a \otimes b$ and $\beta = 1 \otimes b$. $\abs{\alpha} = 3$ and $\abs{\beta} = 2$. Note $a^2 = 0$ and $b^k$ generates $H^{2k}(\CP\infty)$. 
	
	Thus 

	\[\alpha \smile \alpha = (-1)^{\abs{b}\abs{a}}\alpha^2 \otimes b^2 = 0 \otimes b^2 = 0\]

	\[\alpha \smile \beta = (-1)^{\abs{b}\abs{1}} a \otimes b^2 = a \otimes b^2\]

	\[\beta \smile \beta = (-1)^{\abs{b}\abs{1}} 1 \otimes b^2 = 1 \otimes b^2\]

	Thus \[H^*((S^1 \x \CP{\infty}) / (S^1 \x \left\{ x_0 \right\})) = \frac{\Z[\alpha, \beta]}{(\alpha^2)}\]  

	\textbf{Can be extended rather trivially to all coefficients. Basically only part we need to mess with is $H^*(\CP \infty)$ everything else follows with $\R$-Kunneth Formula. } 
\end{proof}

\end{document}